\section{Setup and installation}
\label{sec:setup}

There are two ways to run Learning Mocha: install the shipped APK, which needs nothing but an Android
device, or build from source. The AI assistant additionally needs the small gateway process, and the app
is designed so that everything else keeps working when that process is absent.

\subsection{Installing the shipped APK}

\begin{enumerate}
  \item Copy \code{apk/app-release.apk} from the submission archive to an Android device or emulator
        running \textbf{Android 7.0 (API 24)} or newer.
  \item Install it: \code{adb install -r app-release.apk}, or open the file on the device and allow
        installation from unknown sources.
  \item Launch \emph{Learning Mocha}. On first start the app seeds a small ``Getting Started'' branch so
        no screen is empty, and no account, permission grant or network connection is required.
\end{enumerate}

The APK is signed with a self-generated release key. The keystore and its passwords are deliberately not
part of the archive, so the build is reproducible without them --- a source build without
\code{keystore.properties} simply produces an unsigned release APK instead of failing.

\subsection{Building the Android app from source}

\begin{description}[leftmargin=2.2em, style=nextline, font=\normalfont\bfseries]
  \item[Prerequisites] JDK 17 (the JetBrains Runtime bundled with Android Studio is fine), the Android
    SDK with platform 35 and build-tools 34 or newer. No other tooling is needed --- the Gradle wrapper
    downloads Gradle 8.9 itself.
  \item[SDK location] Create \code{local.properties} in the repository root with
    \code{sdk.dir=<path to your Android SDK>}. Android Studio writes this file automatically the first
    time the project is opened.
  \item[JDK selection] Gradle 8.9 cannot run on JDK 23 or newer. If your default \code{java} is newer,
    either point \code{JAVA\_HOME} at a JDK 17 for the session or add
    \code{org.gradle.java.home=<path to a JDK 17>} to your own \code{\textasciitilde/.gradle/gradle.properties}.
    A machine-specific path is deliberately kept out of the committed \code{gradle.properties} so the
    project configures on any machine.
\end{description}

\begin{lstlisting}[language=bash, caption={Building and testing from the repository root (PowerShell).}]
.\gradlew.bat assembleDebug              # debug APK -> app\build\outputs\apk\debug\
.\gradlew.bat installDebug               # install on a connected device or emulator
.\gradlew.bat test                       # JVM unit tests
.\gradlew.bat connectedDebugAndroidTest  # instrumented Room migration tests (needs a device)
.\gradlew.bat assembleRelease            # release APK -> app\build\outputs\apk\release\
\end{lstlisting}

\subsection{Running the AI gateway}

The gateway is a single Node process whose only job is to hold the DeepSeek API key and shape requests.
It requires Node 18 or newer (developed on Node 24).

\begin{lstlisting}[language=bash, caption={Starting the gateway.}]
cd backend
npm install
copy .env.example .env      # then set DEEPSEEK_API_KEY=<your key>
npm run dev                 # listens on http://localhost:8787
\end{lstlisting}

\code{.env} carries three settings: \code{DEEPSEEK\_API\_KEY} (required), \code{MODEL} (default
\code{deepseek-chat}) and \code{PORT} (default \code{8787}). The file is git-ignored and never packaged,
so the key exists in exactly one place and never inside the APK.

Pointing the app at the gateway:

\begin{itemize}
  \item \textbf{Emulator} --- the default address \code{http://10.0.2.2:8787/} already resolves to the
        host machine's \code{localhost}, so nothing needs configuring.
  \item \textbf{Physical device} --- open \textbf{Settings $\rightarrow$ AI gateway} in the app and enter
        \code{http://<your computer's LAN IP>:8787/}, then use \textbf{Test connection}. The device and
        the computer must be on the same network.
\end{itemize}

The app polls \code{GET /v1/health} before every request. When the gateway is unreachable the AI screens
show a calm banner (``AI unavailable --- your library still works'') with a retry button, and every other
feature behaves exactly as it does online. Health also reports whether the gateway supports streaming, so
an older gateway degrades to the non-streaming endpoint rather than failing.

\subsection{Rebuilding this report}

\begin{lstlisting}[language=bash, caption={The report is plain \LaTeX{} with no external image dependencies.}]
cd report
pdflatex -interaction=nonstopmode report.tex   # run twice so the table of contents settles
\end{lstlisting}

\subsection{Assembling the submission archive}

\code{tools/package-submission.ps1} builds the graded archive from a clean working tree. It stages
\code{README.md}, the full source (with \code{.git} history and without build output), the release APK,
this PDF and the demo-video link into \code{dist/24125006\_24125009.zip}. The script deliberately fails
rather than continuing if a secret --- \code{backend/.env}, the keystore, \code{keystore.properties} or
anything that matches a DeepSeek key pattern --- would end up inside the archive.
